\chapter{Recapitulare șiruri}


\subsection*{Exerciții}
\index{șiruri!convergente} \index{șiruri!limite}

1. Calculați limitele șirurilor cu termenul general:
\begin{multicols}{2}
  \begin{enumerate}[(a)]
  \item $ a_n = \dfrac{n^3 + 5n^2 + 1}{6n^3 + n + 4} $;
  \item $ a_n = \dfrac{n^2 + 2n + 5}{5n^3 - 1} $;
  \item $ a_n = \dfrac{\sqrt{n^2 + 1} + 4\sqrt{n^2 + n + 1}}{n} $;
  \item $ a_n = \dfrac{1^2 + 2^2 + 3^2 + \dots + n^2}{n^3} $;
  \item $ a_n = \dfrac{1 \cdot 2 + 2 \cdot 3 + \dots + n(n+1)}{2n^3} $;
  \item $ a_n = \dfrac{1}{1 \cdot 3} + \dfrac{1}{3 \cdot 5} + \dots + \dfrac{1}{(2n-1)(2n+1)} $;
  \item $ a_n = \sqrt{n^2 + n + 1} - n $;
  \item $ a_n = \sqrt[3]{n^2 + 1} - \sqrt[3]{n^2 + n} $;
  \item $ a_n = \Big( \dfrac{n+1}{2n} \Big)^{\frac{3n^2}{n^2 + 10}} $;
  \item $ a_n = \ln \dfrac{n^3 + 6n}{n^3 + n^2 + n + 1} $;
  \item $ a_n = \dfrac{1}{n^2 + 1} + \dfrac{2}{n^2 + 2} + \dots + \dfrac{n}{n^2 + n} $;
  \item $ a_n = \dfrac{1}{2} + \dfrac{1}{3} + \dfrac{1}{4} + \dots + \dfrac{1}{n} $;
  \item $ a_n = \Big( \dfrac{2n + 1}{2n - 1} \Big)^{n - \sqrt{n}} $;
  \item $ a_n = \dfrac{\ln n}{n} $;
  \item $ a_n = \dfrac{\ln 2 + \ln 3 + \dots + \ln n}{n + 1} $;
  \item $ a_n = \Big( \dfrac{n + 5}{n + 2} \Big)^n $;
  \item $ a_n = \Big( n + \dfrac{1}{n} \Big) \cdot \ln \Big( 1 + \dfrac{n}{n^2 + 1} \Big) $;
  \item $ a_n = \sqrt[n]{n!} $;
  \item $ a_n = \dfrac{n}{\sqrt[n]{n!}} $;
  \item $ a_n = \dfrac{1}{4} + \dfrac{1}{16} + \dfrac{1}{64} + \dots + \dfrac{1}{4^n} $.
  \end{enumerate}
\end{multicols}

\vspace{1cm}

2. Să se determine parametrii $ a, b \in \RR $ astfel încît:
\begin{enumerate}[(a)]
\item $ \ds\lim_{n \to \infty} \Big( \dfrac{n^2 + n + 1}{2n + 1} - an - b \Big) = 1 $;
\item $ \ds\lim_{n \to \infty} \Big( \dfrac{n^2}{n + a} - n \Big) = 1 $;
\item $ \ds\lim_{n \to \infty}\big( a\sqrt{n + 3} + \sqrt{n + 1} \big) = 0 $.
\end{enumerate}

\vspace{1cm}

3. Precizați valoarea de adevăr a propozițiilor de mai jos. În caz de adevăr,
justificați, în caz de falsitate, dați un contraexemplu:
\begin{enumerate}[(a)]
\item Orice șir monoton este mărginit.
\item Orice șir mărginit este monoton.
\item Orice șir convergent este monoton.
\item Orice subșir al unui șir monoton este monoton.
\item Există șiruri monoton crescătoare care au cel puțin un subșir monoton descrescător.
\item Suma a două șiruri monotone este un șir monoton.
\item Diferența a două șiruri monoton crescătoare este un șir monoton crescător.
\item Suma a două șiruri nemărginite este un șir nemărginit.
\item Orice șir divergent este nemărginit.
\item Orice șir nemărginit are limita $ \pm \infty $.
\item Produsul a două șiruri care nu au limită este un șir care nu are limită.
\item Dacă două șiruri sînt convergente, atunci raportul lor este un șir convergent.
\item Dacă pătratul unui șir este convergent, atunci șirul inițial este convergent.
\item Dacă un șir convergent are toți termenii diferiți de zero, atunci limita sa
  este diferită de zero.
\end{enumerate}



%%% Local Variables:
%%% mode: latex
%%% TeX-master: "../bcd-18-19"
%%% End:
